\begin{figure}[tb]
\centering
\begin{tikzpicture}[annotation text]
\node[draw,rounded corners,derived light region,minimum width=4.5cm,minimum height=2.2cm,align=center] (exact) at (0,0)
 {exact\\\([A,B]=0\)\\common eigenvectors\\zero off-diagonal energy};
\node[draw,rounded corners,highlighted light region,minimum width=4.5cm,minimum height=2.2cm,align=center] (approx) at (6.0,0)
 {approximate\\\([A,B]\neq0\)\\weighted Jacobi rotations\\minimum residual coupling};
\draw[coordinate axis,strong line] (exact)--node[above]{generic machine parameters}(approx);
\node[below=8mm of exact] {\(\begin{bmatrix}a&0&0\\0&b&0\\0&0&c\end{bmatrix}\)};
\node[below=8mm of approx] {\(\begin{bmatrix}a&\epsilon&0\\\epsilon&b&\eta\\0&\eta&c\end{bmatrix}\)};
\end{tikzpicture}
\caption{Exact versus approximate simultaneous diagonalization. Residual
off-diagonal terms quantify the compromise rather than disappearing by name.}
\label{fig:joint-diagonalization}
\end{figure}
